\hypertarget{syntax-ergonomics-f-strings-type-annotations}{%
\chapter{SYNTAX ERGONOMICS: F-STRINGS \& TYPE
ANNOTATIONS}\label{syntax-ergonomics-f-strings-type-annotations}}

\hypertarget{pep-498-formatted-string-literals-f-strings}{%
\subsection{12.1 PEP 498: Formatted String Literals
(f-strings)}\label{pep-498-formatted-string-literals-f-strings}}

Introduced in Python 3.6, \textbf{PEP 498} added formatted string
literals (f-strings). These are significantly faster than legacy
formatting methods like \passthrough{\lstinline!\%!} or
\passthrough{\lstinline!.format()!}: * \textbf{Legacy Overheads}:
\passthrough{\lstinline!.format()!} requires a lookup attribute call
(\passthrough{\lstinline!\_\_getattribute\_\_!}) and executes a full
function call frame. * \textbf{f-string Compilation}: The compiler
evaluates f-string expressions at compile time and emits optimized
bytecodes: - \passthrough{\lstinline!FORMAT\_VALUE!}: Formats a single
value on the stack. - \passthrough{\lstinline!BUILD\_STRING!}:
Concatenates values directly in C, bypassing function call overhead.

\begin{lstlisting}[language=Python]
import dis

def legacy_format(name, age):
    return "Name: {}, Age: {}".format(name, age)

def fstring_format(name, age):
    return f"Name: {name}, Age: {age}"

print("Bytecode for f-string:")
dis.dis(fstring_format)
\end{lstlisting}

\begin{lstlisting}
  2           0 LOAD_CONST                1 ('Name: ')
              2 LOAD_FAST                0 (name)
              4 FORMAT_VALUE             0
              6 LOAD_CONST                2 (', Age: ')
              8 LOAD_FAST                1 (age)
             10 FORMAT_VALUE             0
             12 BUILD_STRING             4
             14 RETURN_VALUE
\end{lstlisting}

The f-string compiles to a flat sequence of
\passthrough{\lstinline!FORMAT\_VALUE!} and
\passthrough{\lstinline!BUILD\_STRING!} instructions, achieving
performance comparable to manual string concatenation.

\hypertarget{pep-526-variable-function-type-annotations}{%
\subsection{12.2 PEP 526: Variable \& Function Type
Annotations}\label{pep-526-variable-function-type-annotations}}

\begin{itemize}
\tightlist
\item
  \textbf{The Static Boundary}: Annotations (introduced for functions in
  PEP 3107 and variables in PEP 526) are static metadata. They are not
  checked at runtime and do not impact execution speed.
\item
  \textbf{The Annotation Store}: When a module is loaded, CPython
  compiles annotations and stores them in the class or module namespace
  under the \passthrough{\lstinline!\_\_annotations\_\_!} dictionary.
\item
  \textbf{Inspection}: Type checkers (like Mypy) read this metadata
  statically, while runtime libraries (like Pydantic) use the
  \passthrough{\lstinline!\_\_annotations\_\_!} dictionary or
  \passthrough{\lstinline!inspect.get\_type\_hints()!} to enforce types
  dynamically.
\end{itemize}

\begin{center}\rule{0.5\linewidth}{0.5pt}\end{center}
